\section{Problema e Objetivo}

O problema é a escassez de mão de obra na colheita. A Tevel associa essa restrição à dificuldade de colher no momento adequado e propõe a automação para reduzir a dependência de trabalhadores sazonais~\cite{tevel2026,unifrutti2023tevel}. O desafio é adequar a capacidade de coleta à janela de maturação: atrasos comprometem o objetivo agronômico, enquanto a colheita indiscriminada elimina a seletividade.

Neste estudo, a quantidade de frutos comercializáveis colhidos com qualidade constitui o critério de desempenho operacional. Movimentos executados ou frutos detectados, isoladamente, não representam a efetividade da colheita. Segundo a Unifrutti, a operação em seus pomares realizou colheita seletiva, depositou os frutos em recipientes e minimizou contusões. O comunicado também descreve o acesso à quantidade colhida e a dados de peso, classificação por cor, maturação, diâmetro, horário e geolocalização dos frutos, permitindo caracterizar a produção e sua qualidade~\cite{unifrutti2023tevel}.